\section{Related Works}
\label{sec:related_works}

\subsection{Vision-Language-Action Models}

Recent efforts~\cite{li2024llara,yuan2024robopoint,duan2024aha,niu2024llarva} have adapted vision-language models (VLMs) for action-centric tasks by post-training on curated instruction-following data. For example, RoboPoint~\cite{yuan2024robopoint} and LLARVA~\cite{niu2024llarva} leverage point and visual trajectory into textual prompts to augment LLMs with spatial-action understanding ability. AHA~\cite{duan2024aha} enhances failure detection ability in robotic manipulation by formulating it as a free-form question-answering task, training on synthetic failure data generated by perturbing successful trajectories. Although effective in specific domains, these approaches depend on sophisticatedly curated data and struggle to generalize beyond their training distributions. To improve scalability, recent vision-language-action (VLA) models~\cite{kim24openvla,zheng2024tracevla,szot2024multimodal,bjorck2025gr00t,li2025hamster,yang2025magma,brohan2022rt} adopt large-scale robot datasets (e.g., Open X-Embodiment Dataset~\cite{o2024open} or DROID~\cite{khazatsky2024droid}) to train models directly on diverse demonstrations. OpenVLA~\cite{kim24openvla} learns from pre-trained VLMs with robot trajectories for generalist action execution, while TraceVLA~\cite{zheng2024tracevla} and HAMSTER~\cite{li2025hamster} enhance spatial-action awareness by incorporating visual traces. However, these models predict actions directly from vision and language inputs, often bypassing structured planning or intermediate reasoning. As a result, their capability to handle complex instructions, long-horizon goals, or out-of-distribution scenarios remains limited.


\subsection{Reasoning in Vision-Language-(Action) Models}

Chain-of-thought (CoT) prompting~\cite{wei2022chain,wang2024chain,yeo2025demystifying} has significantly improved the multi-step reasoning ability of LLMs across math, coding, and question-answering tasks. Motivated by these advances, recent works extend reasoning capabilities to vision-language-action (VLA) models for embodied tasks.
ECoT~\cite{zawalski2024robotic} synthesizes intermediate subgoals via prompting and applies supervised fine-tuning to teach VLAs to reason before acting. RAD~\cite{clark2025action} leverages action-free human videos to curate reasoning traces by prompting off-the-shelf LLMs and learn to map reasoning to real actions using robot data. On the other hand, CoT-VLA~\cite{zhao2025cot} replaces linguistic CoT with visual subgoal frames generated ahead of action prediction. However, they depend on either curated CoT supervision or task-specific video generation, limiting their scalability. Inspired by the recent success of RL-optimized reasoning models~\cite{shao2024deepseekmath,guo2025deepseek}, several approaches~\cite{feng2025video,nvidia2025cosmosreason1physicalcommonsense,tan2025reason,liu2025visual} adopt GRPO~\cite{shao2024deepseekmath} optimization to guide CoT generation in vision-language tasks using verifiable rewards. However, their QA-formatted rewards cannot fully support long-horizon planning or establish grounding between reasoning and action execution. To unify structured CoT reasoning with embodied decision-making, we introduce \ours{}, which leverages action-aligned reinforcement learning and visual latent planning to connect embodied reasoning with real-world action in VLA tasks.